% Response to the reviewers, IEEE Systems Journal ISJ-RE-26-22639. Written 2026-11-02 (Session 2).
% Built by a script from the reviewer matrix's verbatim sentence splits. Numbers carry the same markers as the paper.
\documentclass[10pt]{article}
\usepackage[letterpaper,margin=1in]{geometry}
\usepackage{amsmath,amsfonts}
\usepackage{xcolor}
\usepackage{enumitem}
\setlist{nosep,leftmargin=1.4em}
\setlength{\parskip}{4pt}
\setlength{\parindent}{0pt}
\newcommand{\syn}[1]{\textcolor{red}{#1}}            % placeholder data, replaced by the campaign
\newcommand{\rev}[1]{\textcolor{blue}{#1}}           % real result awaiting Chris's review
\newcommand{\open}[1]{\textcolor{magenta}{[#1]}}     % protocol setting not yet decided
% For submission: redefine all three as {#1} (and \open as {}) after the placeholders are gone.
\newcommand{\resp}{\par\textbf{Response.}~}
\newcommand{\chg}{\par\textbf{Changes.}~}
\begin{document}

\begin{center}
{\large\bfseries Response to the Associate Editor and Reviewers}\\[3pt]
Manuscript ISJ-RE-26-22639, ``Adaptive Navigation of Multirobot Systems to Critical Points in 2D Vector Fields in Simulation and Experiment''\\[2pt]
C. Waight and C. A. Kitts
\end{center}

Dear Associate Editor,

We thank you and the three reviewers for their careful reading. The revision adds an error bound for the estimate with an existence test and a covariance (Section~II-B), a proposition for the three-robot minimum (Section~II-C), and a closed-loop analysis with actuator lag, sampled updates, friction, and nearly singular estimates (Section~III-D). Simulation now uses a robot model with Coulomb friction, scores simulation and hardware with one set of metrics (Section~IV-A), and adds nonlinear fields with several critical points (Section~IV-C). Section~VI-B replaces the drift comparison with matched baselines that share the proposed laws, gains, and robot model. \syn{All hardware experiments were repeated in a new campaign designed around the reviewers' requests. All hardware results in the revised manuscript come from this campaign and supersede the earlier results.}

While re-analyzing the original logs we found that the vortex rows of the submitted Tables~IV and~V and the two vortex convergence figures came from four-robot runs. The revision reports three-robot runs only. Each response quotes the comment, split into lettered sentences, and gives the location of each change. Section, equation, table, and figure numbers refer to the revised manuscript. The submitted Tables~IV and~V correspond to the hardware rows of the revised Table~II and to the revised Table~V.


\section*{Associate Editor}

\begin{quote}\itshape
\textbf{(a)}~To sum up, three reviewers have examined the manuscript and found the topic meaningful, but substantial technical revisions are required.
\textbf{(b)}~Please add nonlinear field benchmarks
\textbf{(c)}~and an estimation error analysis that accounts for field curvature, measurement noise, conditioning, and Jacobian invertibility.
\textbf{(d)}~The stability discussion should be strengthened for changing estimates, actuator lag, sampled updates, and nearly singular cases.
\textbf{(e)}~Please also include matched fixed radius baselines,
\textbf{(f)}~repeated trials,
\textbf{(g)}~and clearly defined convergence, localization, radial tracking, control effort, and failure metrics.
\textbf{(h)}~In addition, formalize the three robot minimum,
\textbf{(i)}~clarify the centralized and distributed implementation and the practical scope of the results,
\textbf{(j)}~distinguish the prior framework components from the new contributions,
\textbf{(k)}~and discuss the limitations of the static laboratory setup.
\textbf{(l)}~Please revise the manuscript according to all reviewer comments.
\end{quote}


\resp (a) We thank the Associate Editor. The responses to (b)--(k) below are the full answers. The reviewer sections refer back to them where they ask for the same work.

(b) \emph{Nonlinear benchmarks.} Section~IV-C adds two nonlinear fields. One is a single cell of the double gyre, with a center, four corner saddles, and a curve where $\det\mathbf{J} = 0$. The other is a saddle whose inflow axis is bent by a cubic term of strength $\chi$, in an irrotational and a rotational form. The section reports four results.
\begin{itemize}
\item The estimate's curvature error grows as the cube of the formation side, while the reconstructed fields keep a floor set by their fixed distortion (Fig.~3, left).
\item The two curved saddles respond differently to the same cubic. In the irrotational form the estimate errors on the two axes cancel in the mean orbit radius, and in the rotational form they shift it (Fig.~3, right).
\item The attraction law is a sampled Newton flow, so the sign of $\det\hat{\mathbf{J}}$ at the start predicts which critical point the cluster reaches. The section reports the gyre outcome by start distance against the oracle, and the curved saddle up to 0.7~m.
\item With lag and friction off, added measurement noise leaves the error on the reconstructed fields nearly unchanged, which separates fixed distortion from temporal noise.
\end{itemize}

(c) \emph{Estimation error analysis.} Section~II-B bounds the estimation error in (9). Field curvature enters through a Hessian bound $M$, formation size through $R$, measurement error through $\bar{\epsilon}$, formation conditioning through $\kappa_0$, and Jacobian invertibility through $\sigma_{\min}(\mathbf{J})$, with Weyl's inequality giving invertibility of $\hat{\mathbf{J}}$. The same bound gives a run-time existence test that certifies a true critical point near the estimate. Equation~(10) gives the covariance of the estimate under measurement noise.

(d) \emph{Stability.} Section~III-D models the changing estimate as $\hat{\mathbf{p}}^*[j] = \mathbf{p}^* + \boldsymbol{\delta}[j]$ and analyzes the sampled loop with lag (16). It derives the gain condition (17), treats friction as a bounded command disturbance and the command limit as a bound on velocity, and states Proposition~2, an ultimate bound (18) under explicit assumptions. Nearly singular estimates are handled as described in the response to Reviewer~3, Comment~3(e). The orbital law's radial error is bounded to first order by (19), and the singular point $r = 0$ is treated at the end of the section.

(e) \emph{Matched baselines.} Section~VI-B and Table~VI drive the same attraction and orbital laws, gains, robot model, and formation with three estimates of the critical point. These are the plane fit, the true critical point (an oracle that isolates the controller), and a magnitude-only estimate that senses only the speeds. The orbital comparison is at fixed commanded radii.

(f) \emph{Repeated trials.} Simulation reports 1000 attraction trials per field (Table~II), 100 orbit trials per field and radius (Table~III), and repeated baseline trials (Table~VI). On hardware, \syn{each attraction start and each orbit radius is repeated (Section~V-B), and Table~V separates between-trial from within-run variation}.

(g) \emph{Metrics.} Section~IV-A defines the metrics once for simulation and hardware. These are convergence at thresholds of 0.03, 0.04, and 0.05~m with a failure criterion, convergence time, localization error $\|\hat{\mathbf{p}}^* - \mathbf{p}^*\|$, radial tracking error with a 10~s exclusion and its between-trial and within-run spread, control effort, and classification accuracy.

(h) \emph{Three-robot minimum.} Proposition~1 (Section~II-C) states that three non-collinear robots are necessary and sufficient to locate the critical point of every affine field with $\det\mathbf{J} \neq 0$ from instantaneous measurements.

(i) \emph{Centralized and distributed implementation, practical scope.} The abstract, Section~I, Section~III, and Limitation Five of Section~VI-C state that sensing is spatially distributed while one central computer forms the estimate and computes the commands from OptiTrack positions. The abstract and Section~VII state that the results hold under lab conditions, where the printed fields do not advect the robots.

(j) \emph{Prior framework and new contributions.} Section~III states that the three-layer architecture follows [5], [8] and the cluster space controller follows [42], and that the architecture is not a contribution of this work. The caption of Fig.~2 marks the feature estimation and navigation controller blocks as new. Section~I restates the four contributions.

(k) \emph{Static laboratory setup.} Limitation Two of Section~VI-C names four features of a real flow that the printed maps omit and the effect of each on the method. These are advection of the robots, drifting critical points, communication dropouts, and sensor latency.

(l) Responses to every reviewer comment follow.

\section*{Reviewer 1}
We thank Reviewer~1 for a detailed and technically precise review.

\subsection*{Comment 1.1}

\begin{quote}\itshape
\textbf{(a)}~The scope of the estimation result needs to be clarified.
\textbf{(b)}~Equation (6) exactly recovers the critical point of the fitted affine field, but this point need not correspond to a nearby zero of the true nonlinear field.
\textbf{(c)}~Please establish conditions under which this inference is valid
\textbf{(d)}~and quantify the estimation error in terms of field curvature, formation size, measurement errors, and Jacobian invertibility.
\textbf{(e)}~Validation on nonlinear fields, including fields with multiple critical points, would substantially strengthen the paper.
\end{quote}


\resp (a, b) We agree. The estimate is the zero of the fitted affine field. Section~II-B now relates it to a zero of the true field through the exact relation (7), $\hat{\mathbf{p}}^* - \mathbf{p}^* = -\hat{\mathbf{J}}^{-1}\hat{\mathbf{v}}(\mathbf{p}^*)$. The estimate is therefore displaced by the model error at the true critical point, amplified by $\hat{\mathbf{J}}^{-1}$.

(c, d) Equation~(9) bounds that displacement in terms of field curvature ($M$), formation size ($R$), measurement error ($\bar{\epsilon}$), formation conditioning ($\kappa_0$), and Jacobian invertibility ($\sigma_{\min}(\mathbf{J})$). The inference is valid where the bound's denominator is positive. The existence test that follows (9) certifies, from the current fit and the bounds $M$ and $\bar{\epsilon}$, that a true critical point lies within a stated distance of the estimate. See also the response to the Associate Editor, (c).

(e) Section~IV-C adds nonlinear fields, including a gyre cell with a center and four saddles. The response to the Associate Editor, (b), summarizes the results. Limitation Four of Section~VI-C states which nonlinear regimes remain untested.

\subsection*{Comment 1.2}

\begin{quote}\itshape
\textbf{(a)}~The stability analysis considers an ideal centroid model with a stationary critical point estimate.
\textbf{(b)}~In implementation, the estimate changes with robot positions and sensor measurements, while actuator lag, saturation, stiction, and sampled control introduce additional dynamics.
\textbf{(c)}~The statement in Section VI-A that estimation errors affect accuracy but preserve stability therefore requires justification.
\textbf{(d)}~Please provide a closed-loop robustness analysis with explicit assumptions
\textbf{(e)}~and address the singularity of the orbital law at (r=0).
\end{quote}


\resp (a, b, c) We agree that the submitted Section~VI-A claimed more than the analysis supported. That sentence was removed. The last paragraph of Section~VI-A now states that the estimation error enters the closed loop as an input. Under the conditions of Proposition~2 it does not destabilize the attraction law, and the guarantee lapses as $\hat{\mathbf{J}}$ approaches singularity.

(d) Section~III-D gives the closed-loop robustness analysis with its assumptions stated explicitly.
\begin{itemize}
\item \emph{Changing estimate.} The estimate is $\mathbf{p}^* + \boldsymbol{\delta}[j]$, where $\boldsymbol{\delta}$ depends on the cluster position and on the measurements.
\item \emph{Lag and sampling.} The model is the sampled first-order lag of Section~III-A, which gives the loop (16) and the exact Jury condition (17) on $k\Delta t$.
\item \emph{Friction.} Coulomb friction enters as a command disturbance of norm at most $v_c$.
\item \emph{Saturation.} The command limit bounds the velocity. Proposition~2 assumes that the command near rest is unsaturated.
\end{itemize}
Proposition~2 then bounds the final error by (18), using the error bound (9) for $\boldsymbol{\delta}$. For the orbital law, (19) gives the radial error dynamics to first order and shows that a tangential estimate error is amplified by $k_t/(k_r r_d)$.

(e) The ideal radial dynamics keep $r(t) > 0$, so the singular point $r = 0$ of (12) is not reached. A noisy estimate can still land on the centroid. Below $r = 10^{-6}$~m the implementation fixes the radial unit vector, so the command stays defined, and its tangential step carries the cluster off the singular point on the next cycle (end of Section~III-D).

\subsection*{Comment 1.3}

\begin{quote}\itshape
\textbf{(a)}~The conditioning and noise analysis needs greater rigor.
\textbf{(b)}~Equation (6) involves errors in both the estimated Jacobian and the offset vector, which are obtained from the same measurements and may be correlated.
\textbf{(c)}~Please derive an appropriate perturbation bound or covariance approximation
\textbf{(d)}~and validate it numerically.
\textbf{(e)}~The discussion should also distinguish Jacobian singularity from non-hyperbolicity: a center with purely imaginary, nonzero eigenvalues is non-hyperbolic despite having a nonsingular Jacobian.
\end{quote}


\resp (a, b) The errors in $\hat{\mathbf{J}}$ and $\hat{\mathbf{h}}$ are indeed correlated. Section~II-B shows that they act through one combination. The fitted value at the critical point is $\hat{\mathbf{v}}(\mathbf{p}^*) = \sum_i w_i\mathbf{m}_i$, where $w_i$ are the barycentric coordinates of $\mathbf{p}^*$ with respect to the robots and $\mathbf{m}_i$ are the measurements. The estimation error (7) depends on the measurements only through this combination.

(c) Equation~(9) gives the perturbation bound. For independent measurement noise, (10) gives the first-order covariance
\[
\operatorname{Cov}(\hat{\mathbf{p}}^*) \approx \sigma_v^2\|\mathbf{w}\|^2\mathbf{J}^{-1}\mathbf{J}^{-T}.
\] Here $\|\mathbf{w}\|^2 = 1/3$ at the centroid of any triangle.

(d) Propagated through the sampled loop, the covariance predicts the simulated closed-loop mean error to a median ratio of \rev{1.004} (range \rev{0.98} to \rev{1.04}) over \rev{66} noise settings on the linear fields (Section~IV-C).

(e) Section~VI-A now separates the two. A degenerate critical point has $\det\mathbf{J} = 0$ and $\sigma_{\min}(\mathbf{J}) = 0$, and the bound gives no guarantee there. A center is non-hyperbolic yet has $\sigma_{\min}(\mathbf{J}) > 0$, so (9) locates it as well as any other critical point, and only its type is fragile. The classifier of Section~IV-D therefore treats the center and node boundaries with a stated band.

\subsection*{Comment 1.4}

\begin{quote}\itshape
\textbf{(a)}~The orbital results indicate substantial deviations from the commanded radii, particularly for small radii.
\textbf{(b)}~For example, Table V reports a mean radial error of 0.246 m for a commanded radius of 0.01 m in the vortex field.
\textbf{(c)}~Please define the practically achievable radius range,
\textbf{(d)}~quantify its dependence on actuator dynamics and controller gains,
\textbf{(e)}~and report repeated trials with consistent transient exclusion criteria for simulation and hardware.
\end{quote}


\resp (a, b) We agree. The revision no longer presents small-radius orbits as accurate tracking. Section~V-D, the abstract, and Section~VII describe them as bounded orbits that do not track the commanded radius.

(c) Section~IV-A defines a radius as practically achievable when the magnitude of its radial bias is at most $0.25\,r_d$. Section~IV-B and Section~V-D report the achievable range in simulation and hardware.

(d) Section~IV-B derives the steady orbit radius in closed form (20). The actuator lag delays the rotating tangential command, which turns part of it outward, and friction lowers the orbit speed. In the continuous limit this gives $r^* = (r_d + \sqrt{r_d^2 + 4C})/2$ with $C = \tau V^2/k_r$, so the error approaches $r_{\min} = \sqrt{C}$ at small radii and falls as $C/r_d$ at large ones. The section quantifies how lowering $k_t$, raising $k_r$, or removing the lag reduces $r_{\min}$.

(e) Simulation repeats every radius 100 times on each field (Table~III). Hardware \syn{repeats every radius on both printed maps} (Section~V-B, Table~V). Both exclude the first 10~s.

\subsection*{Comment 1.5}

\begin{quote}\itshape
\textbf{(a)}~The comparative evaluation is insufficient to establish the advantages of the proposed method.
\textbf{(b)}~The direction-based alternatives lack explicit radial error feedback, making drift unsurprising.
\textbf{(c)}~Please include quantitative comparisons for both attraction and orbital navigation under matched sensing and actuator constraints.
\textbf{(d)}~Report convergence time, localization error, radial tracking error, and control effort,
\textbf{(e)}~and clarify how the modified vector-to-scalar baseline is implemented and tuned.
\end{quote}


\resp (a, b) We agree that the vector-sum controller lacks radial feedback, so its drift said little about the proposed method. The vector-sum comparison was removed, together with its figure.

(c, d) Section~VI-B and Table~VI replace it with matched baselines for attraction and orbit. The plane fit, the true critical point, and a magnitude-only estimate drive the same attraction and orbital laws, with the same gains, robot model, formation, and sensed readings. The magnitude-only estimate uses no separate tuning. Table~VI reports success, mean final error, convergence time, control effort, orbit RMSE at three commanded radii, and the estimate's localization error, over repeated trials.

(e) The submitted Section~VI-B described a modified vector-to-scalar primitive, but the comparison figure showed only the vector-sum primitive and the proposed controller, and no results for the vector-to-scalar primitive were reported. Rather than specify that primitive further, we replaced the comparison with the matched baselines above. The magnitude-only baseline is the matched counterpart that senses a scalar reduction of the flow, and Section~VI-B gives its law and gains.

\subsection*{Comment 1.6}

\begin{quote}\itshape
\textbf{(a)}~The experimental statistics and noise model need clearer definitions.
\textbf{(b)}~The reported mean-position bias is not the same as the mean Euclidean localization error across trials.
\textbf{(c)}~Please define the convergence threshold, precision metric, and failure criteria explicitly,
\textbf{(d)}~and report per-trial error distributions.
\textbf{(e)}~In addition, clarify whether the reconstructed MLP fields represent fixed spatial distortion or temporally varying measurement noise;
\textbf{(f)}~these effects should be evaluated separately.
\end{quote}


\resp (a, b) We agree. Table~II now reports the mean $\pm$ standard deviation of the per-trial error $\|\mathbf{p}_c(T) - \mathbf{p}^*\|$ beside the bias of the mean final position, together with the median and 95th percentile.

(c) Section~IV-A defines the convergence thresholds (0.03, 0.04, and 0.05~m), the precision as the RMS distance of the final positions from their mean, and failure as leaving the map or ending outside the threshold. The same definitions are used for simulation and hardware.

(d) Fig.~6 shows the per-trial final errors in hardware as cumulative distributions \syn{(placeholder data until the campaign)}, and Table~II gives their medians and 95th percentiles for simulation and hardware.

(e, f) The reconstructed MLP fields represent fixed spatial distortion with no temporal noise. Sections~IV and~V-A now say so. The two effects are evaluated separately. In simulation with lag and friction off, adding white measurement noise up to 0.03~m/s leaves the error on the reconstructed fields nearly unchanged, so their error is the distortion (Section~IV-C). On the testbed, a parked robot's reading changes by about \rev{$2\times10^{-5}$}~m/s between cycles, against a fixed calibration error of about 0.02~m/s, so the hardware error is also dominated by distortion (Section~VI-A). The hardware campaign includes runs with the formation parked at fixed points to measure the temporal noise directly \open{campaign block E}.

\subsection*{Comment 1.7}

\begin{quote}\itshape
\textbf{(a)}~Several claims should be aligned with the demonstrated scope.
\textbf{(b)}~The implementation uses a central computer and global motion capture, so ``distributed'' should be distinguished from spatially distributed sensing.
\textbf{(c)}~The three-robot minimum applies to instantaneous identification of an unrestricted affine field and should be stated accordingly.
\textbf{(d)}~Finally, critical point classification is not experimentally evaluated,
\textbf{(e)}~and the printed fields do not physically advect the robots.
\textbf{(f)}~Please validate these additional capabilities or qualify the corresponding claims in the abstract, introduction, and conclusion.
\end{quote}


\resp (a, b) The abstract, Section~I, Section~III, and Limitation Five of Section~VI-C now describe the method as cooperative estimation from spatially distributed measurements, computed on one central computer from OptiTrack positions. Limitation Five notes that onboard relative localization and inter-robot communication were not tested.

(c) Proposition~1 states the minimum for instantaneous identification of every affine field with $\det\mathbf{J} \neq 0$, and Section~I states it the same way.

(d) Classification is now evaluated in simulation. Section~IV-D and Table~IV report the fraction of control cycles on which the type of $\hat{\mathbf{J}}$ matches the true type, with and without measurement noise. A center or star node within a stated angular band of its type boundary counts as that type, and the band is stated with the results. On hardware, the campaign logs the type label on every cycle without feeding it to the controller, and Section~V-C reports its accuracy \open{campaign}. Section~I now says that the eigenvalues indicate the type of the critical point.

(e, f) The abstract, Section~III-A, Limitation Two of Section~VI-C, and Section~VII state that the printed fields do not advect the robots and that the results hold under lab conditions. Limitation Two gives the closed-loop error dynamics when a flow does advect the cluster, and the condition under which they become unstable.

\section*{Reviewer 2}
We thank Reviewer~2 for the review and the specific suggestions.

\subsection*{Comment 2.1}

\begin{quote}\itshape
\textbf{(a)}~The cluster space control framework is adopted from earlier publications.
\textbf{(b)}~Please clearly separate the existing layered cluster space infrastructure reused in this paper.
\textbf{(c)}~Also, please genuinely highlight new contributions: instantaneous distributed Jacobian estimation, minimum robot count proof, attraction/orbital control laws built on Jacobian outputs, and the large scale hardware experimental campaign.
\end{quote}


\resp (a, b) Section~III now states that the three-layer architecture, including the feature estimation block, follows [5], [8], and that the cluster space controller follows [42]. It states that the architecture itself is not a contribution of this work. The caption of Fig.~2 marks the feature estimation and navigation controller blocks as the new parts.

(c) Section~I restates four contributions with ordinals. The first is the instantaneous estimate with its error bound and existence test, and the second is the minimality proof. The third is the attraction and orbital laws with the bounded-tracking result, and the fourth is the hardware validation.

\subsection*{Comment 2.2}

\begin{quote}\itshape
\textbf{(a)}~The theoretical derivation handles robot numbers greater than three via least squares, but simulation and hardware only test the minimal three robot case.
\textbf{(b)}~Please add a brief discussion about expected benefits/trade-offs when deploying larger robot teams.
\end{quote}


\resp (a, b) Limitation Three of Section~VI-C now discusses larger teams. With $n > 3$ robots the least-squares fit averages measurement errors across robots and survives the loss of one robot. Each added robot also adds communication and formation-control load, and a wider footprint asks the affine model to hold over more of the field's curvature, which the bound (9) quantifies through $R$.

\subsection*{Comment 2.3}

\begin{quote}\itshape
\textbf{(a)}~The proof that three non collinear robots form the minimum set for unique estimation is given in plain text prose.
\textbf{(b)}~It is better to formalise this as a short proposition/lemma with clear assumptions.
\end{quote}


\resp (a, b) Section~II-C now states Proposition~1 with its assumptions (an affine field with $\det\mathbf{J} \neq 0$ and instantaneous measurements) and gives its proof of necessity and sufficiency.

\subsection*{Comment 2.4}

\begin{quote}\itshape
\textbf{(a)}~For Equation (7), please show the closed loop error dynamics explicitly, to make exponential convergence directly visible.
\end{quote}


\resp (a) Section~II-D now writes the error $\hat{\mathbf{e}} = \mathbf{p}_c - \hat{\mathbf{p}}^*$ and shows that for a stationary estimate the attraction law gives $\dot{\hat{\mathbf{e}}} = -k\hat{\mathbf{e}}$, so $\hat{\mathbf{e}}(t) = \hat{\mathbf{e}}(0)e^{-kt}$. Section~III-D extends this to the sampled loop with lag.

\subsection*{Comment 2.5}

\begin{quote}\itshape
\textbf{(a)}~All hardware tests use static printed vector fields.
\textbf{(b)}~Please explicitly summarise open practical challenges for real world fluid environments: robot advection by physical flow, time varying / drifting critical points, communication drop outs, and non ideal sensor sampling latency.
\textbf{(c)}~Distinguish results obtained under lab conditions from real world applicability.
\end{quote}


\resp (a, b) Limitation Two of Section~VI-C summarizes the four challenges and the effect of each on the method.
\begin{itemize}
\item \emph{Advection.} A flow advects the robots and changes the error dynamics to $\dot{\mathbf{e}} = (\mathbf{J} - k\mathbf{I})\mathbf{e}$, which is unstable when an eigenvalue of $\mathbf{J}$ has real part above $k$, and actuator lag lowers that limit in rotating flow.
\item \emph{Drifting critical points.} The attraction law trails a moving critical point.
\item \emph{Dropouts.} A communication dropout can leave fewer than three readings, and the fit then has no unique solution.
\item \emph{Latency.} Sensor latency mixes readings from different times into one fit.
\end{itemize}

(c) The abstract and Section~VII state that the results hold under lab conditions, and that a physical flow adds a disturbance the present controller does not compensate for.

\section*{Reviewer 3}
We thank Reviewer~3 for the review and its focus on what the evidence supports.

\subsection*{Comment 3.1}

\begin{quote}\itshape
\textbf{(a)}~The main advance seems to be estimating critical-point coordinates rather than providing directional guidance alone.
\textbf{(b)}~In Section II, however, the construction consists of two affine fits followed by Jacobian inversion, and the three-robot minimum is specific to that reconstruction problem.
\textbf{(c)}~Please explain the methodological novelty beyond this construction
\textbf{(d)}~and state the assumptions behind the minimum requirement more precisely.
\textbf{(e)}~The centralized implementation in Section V-A also needs to be distinguished from distributed computation, even though the measurements are spatially distributed.
\end{quote}


\resp (a, b, c) We agree that the construction is two affine fits and a matrix inversion, and Section~I now says so. Its novelty lies in what is built on that construction. The method closes a feedback loop on the joint zero of the two fits at every cycle, bounds how far that zero lies from a true critical point and certifies one nearby (Section~II-B), and bounds how far the cluster settles from it under lag, sampling, and friction (Section~III-D). It then demonstrates both laws in hardware.

(d) Proposition~1 states the assumptions behind the minimum (see the response to the Associate Editor, (h)).

(e) See the response to the Associate Editor, (i).

\subsection*{Comment 3.2}

\begin{quote}\itshape
\textbf{(a)}~The six canonical fields provide a particularly favorable test because they are affine.
\textbf{(b)}~Section IV-B explicitly notes that exact critical-point recovery makes the controller trajectories independent of field type.
\textbf{(c)}~These results therefore do not establish how reliably the method extrapolates from local measurements in nonlinear environments.
\textbf{(d)}~Please add nonlinear benchmarks with different field curvatures
\textbf{(e)}~and starting distances,
\textbf{(f)}~and examine how formation size affects localization accuracy.
\textbf{(g)}~This would test the proposed estimation capability more directly than adding further trials on affine fields.
\end{quote}


\resp (a, b, c) We agree. Section~IV-B now states that, because the six canonical fields are affine, the estimate is exact on them, so they do not show how it extrapolates to a nonlinear field.

(d, e, f) Section~IV-C adds the requested benchmarks. Field curvature is varied through the gyre cell and the cubic coefficient $\chi$ of the curved saddle. Starting distance is varied in both fields, and the results show where the gyre's $\det\mathbf{J} = 0$ boundary sends the cluster to a neighboring saddle. Formation side is swept well below and above the hardware value (Fig.~3, left), and the localization error grows as the cube of the side on curved fields. The response to the Associate Editor, (b), summarizes the results.

(g) We agree, and these benchmarks test the estimate itself. On every linear field the plane fit and the oracle give identical results, so the differences reported on the nonlinear fields are estimation error.

\subsection*{Comment 3.3}

\begin{quote}\itshape
\textbf{(a)}~The stability discussion makes a stronger claim than the analysis supports.
\textbf{(b)}~Section II-D limits the attraction convergence result to a stationary estimate, but Section VI-A states that estimation and conditioning errors preserve stability.
\textbf{(c)}~An estimate that changes with noisy measurements and robot positions cannot simply be treated as a fixed displacement of the equilibrium.
\textbf{(d)}~Please establish conditions for bounded tracking errors with actuator lag and sampled updates,
\textbf{(e)}~and explain how nearly singular Jacobian estimates are handled.
\end{quote}


\resp (a, b, c) We agree, and the unsupported sentence in Section~VI-A was removed. Section~III-D treats the estimate as $\mathbf{p}^* + \boldsymbol{\delta}[j]$, with a position-dependent part bounded through (9), rather than as a fixed displacement of the equilibrium.

(d) Proposition~2 gives conditions for bounded tracking error with actuator lag and sampled updates. These are the gain condition (17), a held formation, an unsaturated command near rest, a Hessian bound on the region visited, $\underline{\sigma}_J > 0$, and $\gamma\eta < 1$. Under them the final error obeys the ultimate bound (18). See also the response to Reviewer~1, Comment~2.

(e) The fit returns no estimate when $|\det\hat{\mathbf{J}}| < 10^{-10}$, and the controller then \open{holds the last valid estimate}. Otherwise the memoryless estimator replaces a spurious estimate on the next cycle, and the command limit caps the centroid speed whatever the estimate. No conditioning gate on $\sigma_{\min}(\hat{\mathbf{J}})$ was applied, and Section~III-D says so. Proposition~2 gives no guarantee as $\hat{\mathbf{J}}$ approaches singularity. Such cases are detectable at run time through $\sigma_{\min}(\hat{\mathbf{J}})$ and the existence test (Section~VI-A).

\subsection*{Comment 3.4}

\begin{quote}\itshape
\textbf{(a)}~The comparison in Section VI-B is not yet convincing as evidence of algorithmic superiority.
\textbf{(b)}~Neither the vector-sum controller nor the modified vector-to-scalar controller explicitly regulates a commanded radius, so their outward drift mainly demonstrates the absence of radial feedback.
\textbf{(c)}~Please include a baseline designed for the same fixed-radius task under comparable sensing assumptions
\textbf{(d)}~and report quantitative results over repeated trials.
\textbf{(e)}~A separate comparison using the true critical-point location would help distinguish estimation errors from limitations of the orbital controller itself.
\end{quote}


\resp (a, b) We agree that the drift of controllers without radial feedback mainly shows that absence. The vector-sum comparison was removed.

(c, d) Section~VI-B and Table~VI add a baseline for the same fixed-radius task under comparable sensing. A magnitude-only estimate of the critical point, formed from the same three readings, drives the same orbital law with the same gains, robot model, and formation. The table reports it over repeated trials, in attraction and in orbits at three commanded radii, beside the plane fit.

(e) Table~VI also drives both laws with the true critical point. On the affine fields the plane fit matches this oracle in every metric, so the remaining error belongs to the controller. On the reconstructed fields the difference between the two isolates the estimation error. Sections~IV-B and~IV-C use the same oracle on the canonical and nonlinear fields.

\subsection*{Comment 3.5}

\begin{quote}\itshape
\textbf{(a)}~The accuracy claims need more careful interpretation.
\textbf{(b)}~Table IV reports the displacement of the average final position, which is not the same as the average distance error across individual trials.
\textbf{(c)}~Table V also gives a radial error of 0.246 m for a commanded radius of 0.01 m, so bounded orbital motion should not be presented as accurate radius tracking at that setting.
\textbf{(d)}~Please report per-trial distance errors,
\textbf{(e)}~define a numerical convergence threshold,
\textbf{(f)}~and repeat the orbital conditions to distinguish between-trial variability from fluctuations within one trajectory.
\end{quote}


\resp (a, b) We agree. Table~II now reports the per-trial mean error $\pm$ standard deviation, the median, and the 95th percentile beside the bias of the mean final position, and Section~V-C explains why the bias is smaller.

(c) Small-radius orbits are now described as bounded orbits that do not track the commanded radius. Section~V-D reports the practically achievable range, and Section~IV-B explains, through the closed form (20), why the actuator lag sets a smallest achievable orbit.

(d, e) Per-trial distances appear in Table~II and Fig.~6, and Section~IV-A defines the convergence thresholds numerically.

(f) \syn{Each orbital condition is repeated}, and Table~V reports the between-trial spread of the run means beside the within-run standard deviation (Section~V-D).


Sincerely,\\
Christopher Waight and Christopher A. Kitts

\end{document}
